\begin{remark}[What the locality uses and what it fixes]
\label{rem:c11-local}
The argument uses $sO=0$ exactly, the classical master equation, the anticommutation of $Q$ with $\partial$ of the composite kernel, which rests on $[v_x.,Q]=0$ and $QP\star+P\star Q=e^{-\varepsilon H}$ (Lemma \ref{lem:c11-kernels} (v)), the absence of $\log\varepsilon$ in cycles with one special edge, a property of the cutoff on every edge (Theorem \ref{thm:c11-simplex}), the absence of a term $\varepsilon^{-1}$ in the composite loops of $I$ (Theorem \ref{thm:c11-local} (ii)), and generic momenta at $T=\infty$ for the convergence of the cycles and of the composite loops, which are of order $|p|^{-3}$ (Proposition \ref{prop:c9-composite} (iv)), or any finite $T$ (Theorem \ref{thm:c8-L} (ii)); (i) of Definition \ref{def:c10-hyp} enters through the terms with $QW$ of the proof of Theorem \ref{thm:c10-D1}, which are outside the composite term. For a functional of the fields $A|_V$, $A|_{W_{\mathbb R}^\perp}$, $\lambda$ alone, $sO=0$ is the gauge invariance of its density. Statement (iv) fixes the one-body part of the composite term: a term $t$ units times $v.X_a$ on every letter, added to the right side of \eqref{eq:c11-local}, would replace each sector coefficient $\theta$ by $\theta+t$ in Sections 11.3--11.5 and leave every spectrum unchanged; the operator-level value $\theta=-\tfrac13$ of Proposition \ref{prop:c11-theta-sector} is the value in the normalisation of \eqref{eq:c11-local}, and it agrees with the action-level value of Theorem \ref{thm:c9-theta2}. The steps of the proof are Lemma \ref{lem:c11-Cid} for the composite loops of the action, Lemma \ref{lem:c11-trees} for the special edges on trees and Proposition \ref{prop:c11-typeI} for the cancellation of the brackets $\pm\epsilon_\Xi\{\Xi_x,O\}$ between the two types, forced by Proposition \ref{prop:c11-master}, whose right side contains no bracket; the sign $\epsilon_\Xi$ never needs to be evaluated. The composite term enters the second evaluation of $D_1$ and $P_1$; the components themselves are given by the master formula of Theorem \ref{thm:c11-classes}, which holds along every generator. Statement (iv) is the input of Corollary \ref{cor:c11-D1local}, Propositions \ref{prop:c11-routes} and \ref{prop:c11-scalar} and Sections 11.3--11.5; its consequences are evaluated on the bilinears, where the single chain of Example \ref{ex:c11-bilinear} reproduces the value of the route of the special edges, on the planar scalar sector, and through the spectra of Table \ref{tab:c11-results}, which a common one-body term leaves unchanged.
\end{remark}
